\chapter{Control Flow \& Preprocessor}
\label{ch:control-flow-preprocessor}

\textit{Making decisions, repeating actions, and bending the compiler to your will before it even runs.}

Every program you have written so far has been a straight line---executing from top to bottom, never deviating. Real software does not work this way. A trading engine must react to market signals in microseconds. A game must loop sixty times per second, branching on player input. An operating system must route interrupts to the correct handler without a single wasted cycle. This chapter equips you with the complete control-flow toolkit of C++98/03 and the preprocessor machinery that operates \textit{before} your code ever reaches the compiler.

%% ============================================================
\section{Conditional Branching: \texttt{if}, \texttt{else if}, \texttt{else}}
\label{sec:if-else}

The \texttt{if} statement is the most fundamental decision-making tool in programming. It evaluates a boolean condition and executes a block of code only if the condition is \texttt{true}.

\begin{lstlisting}[language=C++, caption={Basic if-else chain}]
#include <iostream>

int main() {
    int health = 45;

    if (health > 50) {
        std::cout << "You feel fine." << std::endl;
    } else if (health > 20) {
        std::cout << "You are injured. Find a medkit." << std::endl;
    } else {
        std::cout << "Critical warning! Health is dangerously low." << std::endl;
    }

    return 0;
}
\end{lstlisting}

The conditions are evaluated from top to bottom. As soon as one condition evaluates to \texttt{true}, its block executes, and the rest of the chain is completely skipped. The \texttt{else} block is a catch-all that executes only when every preceding condition evaluated to \texttt{false}.

\subsection{The Boolean Nature of Conditions}

In C++, the \texttt{if} condition can be any expression convertible to \texttt{bool}. Integers are truthy if non-zero, and pointers are truthy if non-null:

\begin{lstlisting}[language=C++, caption={Truthiness of non-boolean types}]
int x = 42;
if (x) {
    std::cout << "x is non-zero" << std::endl;  // This executes
}

int* ptr = NULL;
if (ptr) {
    std::cout << "ptr is valid" << std::endl;
} else {
    std::cout << "ptr is null" << std::endl;      // This executes
}
\end{lstlisting}

%% ============================================================
\section{The Ternary Operator \texttt{?:}}
\label{sec:ternary}

When you need to assign a value based on a simple condition, a full \texttt{if/else} block can feel needlessly verbose. The \textbf{ternary operator} condenses the pattern into a single expression.

\begin{lstlisting}[language=C++, caption={Ternary operator for conditional assignment}]
int player_score = 8500;
int high_score = 10000;

// Syntax: condition ? value_if_true : value_if_false;
int new_high_score = (player_score > high_score) ? player_score : high_score;
\end{lstlisting}

\textbf{Warning:} Do not nest ternary operators. Code like \texttt{(age < 18) ? "Minor" : (age < 65) ? "Adult" : "Senior"} is difficult to read. Use \texttt{if/else} instead.

\begin{lstlisting}[language=C++, caption={Ternary operator examples}]
#include <iostream>

int main() {
    int x = 10, y = 5;

    int max = (x > y) ? x : y;
    std::cout << "Max: " << max << std::endl;  // 10

    std::cout << (x % 2 == 0 ? "Even" : "Odd") << std::endl;  // Even

    return 0;
}
\end{lstlisting}

%% ============================================================
\section{\texttt{switch} and \texttt{case}}
\label{sec:switch}

When you have a single integer or character and want to check it against many possible exact values, a \texttt{switch} statement is cleaner---and often faster---than a massive chain of \texttt{else if} statements. The compiler can optimize a \texttt{switch} into a \textbf{jump table}, yielding O(1) dispatch.

\begin{lstlisting}[language=C++, caption={Basic switch statement}]
#include <iostream>

int main() {
    int day = 3;

    switch (day) {
        case 1:
            std::cout << "Monday" << std::endl;
            break;
        case 2:
            std::cout << "Tuesday" << std::endl;
            break;
        case 3:
            std::cout << "Wednesday" << std::endl;
            break;
        default:
            std::cout << "Unknown day" << std::endl;
            break;
    }

    return 0;
}
\end{lstlisting}

\subsection{The \texttt{break} Imperative and Fall-Through}

Without \texttt{break;}, C++ will \textbf{fall through} and execute the code for the next case, even if the value does not match. This historical quirk inherited from C has caused billions of dollars in software bugs.

Sometimes, you actually want cases to fall through---for example, stacking multiple cases together:

\begin{lstlisting}[language=C++, caption={Intentional fall-through in switch}]
char grade = 'B';

switch (grade) {
    case 'A':
    case 'B':
    case 'C':
        std::cout << "You passed!" << std::endl;
        break;
    case 'F':
        std::cout << "Please see the professor." << std::endl;
        break;
    default:
        std::cout << "Invalid grade." << std::endl;
        break;
}
\end{lstlisting}

\subsection{The \texttt{default} Label}

The \texttt{default} label introduces a block that will be jumped to if the condition's value does not equal any of the \texttt{case} labels' values. The condition must be an expression or declaration of integer or enumeration type, or a class type with a conversion function to integer or enumeration type.

%% ============================================================
\section{\texttt{while} and \texttt{do-while} Loops}
\label{sec:while-loops}

Use a \texttt{while} loop when you want to repeat an action but do not know exactly how many times---you only know it should stop when a condition becomes \texttt{false}.

\begin{lstlisting}[language=C++, caption={While loop}]
#include <iostream>

int main() {
    int ammo = 3;

    while (ammo > 0) {
        std::cout << "Bang!" << std::endl;
        ammo--;
    }
    std::cout << "Click. Out of ammo." << std::endl;

    return 0;
}
\end{lstlisting}

A \textbf{\texttt{do-while}} loop is identical, except the condition is checked at the \textit{end} of the loop body, not the beginning. This guarantees the code will run \textbf{at least once}, even if the condition is \texttt{false} from the start.

\begin{lstlisting}[language=C++, caption={Do-while loop for input validation}]
#include <iostream>

int main() {
    int choice;
    do {
        std::cout << "Press 1 to start, 0 to exit: ";
        std::cin >> choice;
    } while (choice != 0 && choice != 1);

    return 0;
}
\end{lstlisting}

\subsection{\texttt{while} vs.\ \texttt{do-while} Semantics}

The distinction is critical. A \texttt{while} loop with a false condition at entry never executes its body. A \texttt{do-while} loop always executes at least once.

Do not forget the semicolon at the end of \texttt{while(condition);} in the \texttt{do-while} construct.

%% ============================================================
\section{\texttt{for} Loops}
\label{sec:for-loops}

When you know exactly how many times you want to iterate, use a \texttt{for} loop. It packs initialization, condition, and increment into a single, clean line.

\begin{lstlisting}[language=C++, caption={Basic for loop}]
for (int i = 0; i < 5; i++) {
    std::cout << "Count: " << i << std::endl;
}
// Note: The variable 'i' dies here. It only exists inside the loop!
\end{lstlisting}

\subsection{Advanced \texttt{for} Loop Patterns}

The \texttt{for} loop header supports multiple variable declarations and multiple iteration expressions:

\begin{lstlisting}[language=C++, caption={Multi-variable for loop}]
for (int a = 0, b = 10, c = 20; (a + b + c < 100); c--, b++, a += c) {
    std::cout << a << " " << b << " " << c << std::endl;
}
\end{lstlisting}

\subsection{Variable Hiding in \texttt{for} Loops}

A variable declared in the \texttt{for} initialization hides any variable of the same name in the enclosing scope. To use an already-declared variable, omit the type declaration.

\subsection{The \texttt{for} Loop as a \texttt{while} Loop}

Every \texttt{for} loop is equivalent to a \texttt{while} loop: \texttt{for(init; cond; incr) \{body;\}} becomes \texttt{init; while(cond) \{body; incr;\}}.

\subsection{Infinite Loops}

An empty \texttt{for} header creates an infinite loop: \texttt{for (;;) \{...\}}. The equivalent \texttt{while} form is \texttt{while (true) \{...\}}.

\subsection{Iterator-Based Loops (C++98)}

Before C++11, iterating over STL containers required explicit iterator declarations:

\begin{lstlisting}[language=C++, caption={Iterator-based loop over std::vector}]
#include <iostream>
#include <vector>
#include <string>

int main() {
    std::vector<std::string> names;
    names.push_back("Albert Einstein");
    names.push_back("Stephen Hawking");

    for (std::vector<std::string>::iterator it = names.begin();
         it != names.end(); ++it) {
        std::cout << *it << std::endl;
    }

    return 0;
}
\end{lstlisting}

%% ============================================================
\section{Loop Control: \texttt{break} and \texttt{continue}}
\label{sec:break-continue}

\begin{itemize}
    \item \textbf{\texttt{break}}: Instantly destroys the loop. Execution jumps to the first line after the loop block.
    \item \textbf{\texttt{continue}}: Instantly skips the rest of the current iteration and jumps back to the loop's condition check.
\end{itemize}

\begin{lstlisting}[language=C++, caption={break and continue}]
for (int i = 1; i <= 10; i++) {
    if (i == 3) {
        continue;  // Skips printing 3
    }
    if (i == 8) {
        break;     // Destroys the loop at 8
    }
    std::cout << i << " ";
}
// Output: 1 2 4 5 6 7
\end{lstlisting}

Because control flow changes are sometimes difficult for humans to reason about, \texttt{break} and \texttt{continue} should be used sparingly. A more straightforward implementation is usually easier to read.

%% ============================================================
\section{The \texttt{goto} Statement}
\label{sec:goto}

C++ supports the \texttt{goto} statement, which jumps execution to an arbitrary label within the current function. \textbf{Never use \texttt{goto} for control flow.} It creates ``spaghetti code'' that is impossible to follow and debug.

\subsection{The One Acceptable Use: Error Cleanup in C-Style Code}

In legacy C-style resource management (before RAII), \texttt{goto} was sometimes used to jump to a single cleanup block:

\begin{lstlisting}[language=C++, caption={goto for error cleanup (C-style pattern)}]
#include <iostream>
#include <cstdlib>
#include <cstdio>

int main() {
    FILE* file = NULL;
    char* buffer = NULL;

    file = fopen("test.txt", "r");
    if (!file) {
        std::cout << "Failed to open file" << std::endl;
        goto cleanup;
    }

    buffer = new char[100];
    if (!buffer) {
        std::cout << "Memory allocation failed" << std::endl;
        goto cleanup;
    }

    // Do work...

cleanup:
    if (buffer) delete[] buffer;
    if (file) fclose(file);

    return 0;
}
\end{lstlisting}

In modern C++, RAII and exceptions make \texttt{goto} entirely unnecessary. See Chapter~9 for details.

%% ============================================================
\section{Nested Control Flow and the Arrow Anti-Pattern}
\label{sec:nested-flow}

You can nest loops inside loops and \texttt{if} statements inside \texttt{if} statements. Be wary of deeply nested control flow. If your code looks like a giant sideways arrow due to many nested blocks, your code is unreadable.

The solutions:
\begin{enumerate}
    \item Use \textbf{early returns}. If a condition fails, \texttt{return} or \texttt{continue} immediately.
    \item Break inner loops out into separate \textbf{functions}.
\end{enumerate}

%% ============================================================
\section{Operator Precedence and Short-Circuit Evaluation}
\label{sec:operator-precedence}

Understanding operator precedence is not optional for professional C++ engineering. Bugs born from incorrect assumptions about evaluation order are among the hardest to diagnose.

\subsection{Short-Circuit Evaluation}

C++ uses \textbf{short-circuit evaluation} in \texttt{\&\&} and \texttt{||}:
\begin{itemize}
    \item \texttt{\&\&}: If the left operand is \texttt{false}, the right operand is \textit{never evaluated}.
    \item \texttt{||}: If the left operand is \texttt{true}, the right operand is \textit{never evaluated}.
\end{itemize}

\begin{lstlisting}[language=C++, caption={Short-circuit evaluation and precedence}]
#include <iostream>
#include <string>

bool True(std::string id) {
    std::cout << "True" << id << std::endl;
    return true;
}

bool False(std::string id) {
    std::cout << "False" << id << std::endl;
    return false;
}

int main() {
    bool result;

    // Equivalent to: False("A") || (False("B") && False("C"))
    result = False("A") || False("B") && False("C");
    // Output: FalseA, FalseB  (C is skipped)

    // Equivalent to: True("A") || (False("B") && False("C"))
    result = True("A") || False("B") && False("C");
    // Output: TrueA  (B and C are both skipped)

    return 0;
}
\end{lstlisting}

\subsection{Unary Operators and Postfix Semantics}

\textbf{Unary operators} act on a single operand and have high precedence. When used in postfix form, the side-effect occurs only after the entire expression is evaluated:

\begin{lstlisting}[language=C++, caption={Prefix vs.\ postfix increment}]
int a = 1;
++a;            // a is now 2
a--;            // a is now 1
int minusa = -a;  // minusa is -1

a = 4;
int c = a++ / 2;  // c = 4/2 = 2, then a becomes 5
int d = ++a / 2;  // a becomes 6, then d = 6/2 = 3
\end{lstlisting}

\textbf{Undefined Behavior Warning}: Never write expressions like \texttt{a = b++ + ++b;}. Modifying a variable twice between sequence points is Undefined Behavior.

\subsection{Arithmetic Precedence}

Arithmetic operators follow mathematical precedence: multiplication and division before addition and subtraction, all with left-to-right associativity.

\begin{lstlisting}[language=C++, caption={Arithmetic precedence}]
int a = 2 + 4 / 2;       // 2 + (4/2) = 4
int b = (3 + 3) / 2;     // (3+3)/2 = 3
int c = 3 + 4 / 2 * 6;   // 3 + ((4/2)*6) = 15
int d = 3 * (3 + 6) / 9; // (3*(3+6))/9 = 3
\end{lstlisting}

\subsection{Logical AND/OR Precedence}

AND (\texttt{\&\&}) binds tighter than OR (\texttt{||}). In professional code, always parenthesize mixed \texttt{\&\&}/\texttt{||} expressions:

\begin{lstlisting}[language=C++, caption={AND before OR}]
bool can_drive = has_domestic_license || (has_foreign_license && num_days <= 60);
\end{lstlisting}

%% ============================================================
\section{Bitwise Operators}
\label{sec:bitwise}

Bitwise operators manipulate individual bits of integer types. They are essential for embedded systems, graphics programming, cryptography, and network protocol implementation.

\subsection{\texttt{|} --- Bitwise OR}

Each bit in the result is \texttt{1} if \textit{either} corresponding bit is \texttt{1}:

\begin{lstlisting}[language=C++, caption={Bitwise OR}]
int a = 5;      // 0101b
int b = 12;     // 1100b
int c = a | b;  // 1101b = 13
std::cout << "a = " << a << ", b = " << b << ", c = " << c << std::endl;
\end{lstlisting}

\subsection{\texttt{\^{}} --- Bitwise XOR}

Each bit in the result is \texttt{1} if the corresponding bits are \textit{different}:

\begin{lstlisting}[language=C++, caption={Bitwise XOR}]
int a = 5;      // 0101b
int b = 9;      // 1001b
int c = a ^ b;  // 1100b = 12
\end{lstlisting}

\textbf{XOR Swap} --- a classic trick that swaps two variables without a temporary:

\begin{lstlisting}[language=C++, caption={XOR swap (demonstration only)}]
void doXORSwap(int& a, int& b) {
    if (&a != &b) {  // Swapping a variable with itself zeros it!
        a ^= b;
        b ^= a;
        a ^= b;
    }
}
\end{lstlisting}

On modern CPUs, the XOR swap is slower than using a temporary variable due to pipeline stalls. Use \texttt{std::swap()} instead.

\subsection{\texttt{\&} --- Bitwise AND}

Each bit in the result is \texttt{1} only if \textit{both} corresponding bits are \texttt{1}:

\begin{lstlisting}[language=C++, caption={Bitwise AND}]
int a = 6;      // 0110b
int b = 10;     // 1010b
int c = a & b;  // 0010b = 2
\end{lstlisting}

\subsection{\texttt{<<} --- Left Shift}

Shifts all bits left by N positions, padding the least significant bits with zeros. Equivalent to multiplying by $2^N$:

\begin{lstlisting}[language=C++, caption={Left shift}]
int a = 1;       // 0001b
int b = a << 1;  // 0010b = 2
int c = a << 4;  // 00010000b = 16
\end{lstlisting}

Left-shifting a signed number so that the sign bit is affected is \textbf{Undefined Behavior}. Shifting by a negative amount or by more bits than the type can hold is also UB.

\subsection{\texttt{>>} --- Right Shift}

Shifts all bits right by N positions. For unsigned types, zeros are shifted in. For signed negative numbers, the behavior is \textbf{implementation-defined}:

\begin{lstlisting}[language=C++, caption={Right shift}]
int a = 8;       // 1000b
int b = a >> 1;  // 0100b = 4

int x = -2;
int y = x >> 1;  // Result depends on the compiler!
\end{lstlisting}

%% ============================================================
\section{Preprocessor Directives: \texttt{\#define} and \texttt{\#include}}
\label{sec:preprocessor}

The \textbf{preprocessor} operates on your source code \textit{before} the compiler sees it. It performs textual substitution, file inclusion, and conditional compilation. Every line beginning with \texttt{\#} is a preprocessor directive.

\subsection{\texttt{\#define} --- Object-Like and Function-Like Macros}

\begin{lstlisting}[language=C++, caption={Basic preprocessor macros}]
#define PI 3.14159
#define MAX_SIZE 100
#define SQUARE(x) ((x) * (x))

#define DEBUG

#ifdef DEBUG
    #define LOG(msg) std::cout << msg << std::endl
#else
    #define LOG(msg)  // Expands to nothing in release builds
#endif

#include <iostream>

int main() {
    std::cout << "PI = " << PI << std::endl;

    int arr[MAX_SIZE];
    std::cout << "Array size: " << sizeof(arr) << std::endl;

    std::cout << "Square of 5: " << SQUARE(5) << std::endl;

    LOG("Debug message");

    return 0;
}
\end{lstlisting}

\textbf{Macro Pitfall}: Macros perform textual substitution. Always parenthesize macro parameters.

%% ============================================================
\section{Conditional Compilation}
\label{sec:conditional-compilation}

Conditional compilation directives allow you to include or exclude blocks of code at compile time:

\begin{lstlisting}[language=C++, caption={Platform-specific conditional compilation}]
#include <iostream>

#ifdef _WIN32
    #define OS "Windows"
#elif __APPLE__
    #define OS "macOS"
#elif __linux__
    #define OS "Linux"
#else
    #define OS "Unknown"
#endif

int main() {
    std::cout << "Running on: " << OS << std::endl;

#if defined(DEBUG)
    std::cout << "Debug mode" << std::endl;
#else
    std::cout << "Release mode" << std::endl;
#endif

    return 0;
}
\end{lstlisting}

%% ============================================================
\section{Pragma Directives}
\label{sec:pragma}

\texttt{\#pragma} directives provide compiler-specific instructions:

\begin{lstlisting}[language=C++, caption={Pragma directives}]
#include <iostream>

#pragma warning(disable: 4996)  // MSVC

#pragma pack(1)

struct PackedData {
    char a;      // 1 byte
    int b;       // 4 bytes
    double c;    // 8 bytes
};

#pragma pack()   // Restore default packing

int main() {
    std::cout << "Size of PackedData: " << sizeof(PackedData) << std::endl;
    // Without pragma pack: likely 24 (with alignment padding)
    // With pragma pack(1): 13 (tightly packed)

    return 0;
}
\end{lstlisting}

%% ============================================================
\section{Inline Functions vs.\ Macros}
\label{sec:inline-vs-macros}

Macros and inline functions both aim to eliminate function-call overhead, but they are fundamentally different mechanisms.

\begin{lstlisting}[language=C++, caption={Macro vs.\ inline function}]
#include <iostream>

// Macro function --- preprocessor substitution (no type safety)
#define ADD_MACRO(a, b) ((a) + (b))

// Inline function --- type-safe, debuggable
inline int add_inline(int a, int b) {
    return a + b;
}

int main() {
    std::cout << ADD_MACRO(5, 3) << std::endl;    // 8
    std::cout << add_inline(5, 3) << std::endl;   // 8

    // Macro danger: side effects with increment operators
    int x = 5, y = 3;
    std::cout << ADD_MACRO(x++, y++) << std::endl;

    // Inline function is safer
    x = 5; y = 3;
    std::cout << add_inline(x++, y++) << std::endl;

    return 0;
}
\end{lstlisting}

\subsection{The \texttt{do \{ ... \} while(0)} Macro Idiom}

Multi-statement macros must be wrapped in \texttt{do \{ ... \} while(0)} to behave correctly in all contexts:

\begin{lstlisting}[language=C++, caption={Safe multi-statement macro}]
#define BAD_MACRO(x)  f1(x); f2(x); f3(x);
// Dangerous: if (cond) BAD_MACRO(var);  // Only f1 is inside the if!

#define GOOD_MACRO(x)  do { f1(x); f2(x); f3(x); } while(0)
// Safe: if (cond) GOOD_MACRO(var);  // All three calls protected
\end{lstlisting}

%% ============================================================
\section{Professional Insights: Loop Mechanics In Depth}
\label{sec:loop-mechanics}

\subsection{Variable Declaration in Loop Conditions}

In \texttt{for} and \texttt{while} loops, you may declare a variable in the condition. The variable is in scope until the loop ends and persists through each iteration. This is not permitted in \texttt{do-while} loops.

\subsection{The \texttt{for} Loop Execution Model}

The \texttt{for} loop has three optional clauses: initialization (once), condition (before each iteration), and increment (after each iteration). Any or all may be omitted.

\subsection{Breaking Out of Nested Loops}

\texttt{break} only exits the innermost enclosing loop. To break out of multiple nested loops:
\begin{enumerate}
    \item Extract the nested loops into a function and use \texttt{return}.
    \item Use a flag variable checked by the outer loop.
    \item In extreme low-level code, \texttt{goto} to a label after the outer loop.
\end{enumerate}

%% ============================================================
\section{Professional Insights: Flow Control Keywords Reference}
\label{sec:flow-keywords}

\subsection{\texttt{try}, \texttt{catch}, and \texttt{throw}}

Exception handling keywords provide structured error propagation. They are covered in depth in Chapter~9.

\begin{lstlisting}[language=C++, caption={Exception handling basics}]
#include <iostream>
#include <vector>
#include <stdexcept>

void print_asterisks(int count) {
    if (count < 0) {
        throw std::invalid_argument("count cannot be negative!");
    }
    while (count--) { putchar('*'); }
}

int main() {
    try {
        std::vector<int> v(1000000000);
    } catch (const std::bad_alloc&) {
        std::cout << "failed to allocate memory!" << std::endl;
    } catch (const std::runtime_error& e) {
        std::cout << "runtime error: " << e.what() << std::endl;
    } catch (...) {
        std::cout << "unexpected exception!" << std::endl;
        throw;  // Rethrow to caller
    }

    return 0;
}
\end{lstlisting}

\subsection{\texttt{throw} in Expressions}

A \texttt{throw} expression has type \texttt{void} and can be nested in ternary expressions:

\begin{lstlisting}[language=C++, caption={throw in ternary expression}]
unsigned int predecessor(unsigned int x) {
    return (x > 0) ? (x - 1)
                    : (throw std::invalid_argument("0 has no predecessor"));
}
\end{lstlisting}

%% ============================================================
\section{Forward Reference: C++11/17 Control Flow Enhancements}
\label{sec:forward-ref}

The following features are \textbf{not available in C++98/03} but are included as forward references. They will be covered in the C++11 and C++17 volumes.

\subsection{Range-Based \texttt{for} Loops [C++11]}

C++11 introduced the range-based \texttt{for} loop, which drastically simplifies iteration over containers:

\begin{lstlisting}[language=C++, caption={Range-based for loop [C++11]}]
#include <vector>
#include <iostream>

int main() {
    std::vector<double> prices;
    prices.push_back(19.99);
    prices.push_back(5.50);
    prices.push_back(42.00);

    for (double price : prices) {
        std::cout << "Price: $" << price << std::endl;
    }

    return 0;
}
\end{lstlisting}

\subsection{\texttt{if} and \texttt{switch} with Initializers [C++17]}

C++17 allows initialization statements directly inside \texttt{if} and \texttt{switch} conditions, keeping variable scope minimal.

\subsection{\texttt{[[fallthrough]]} Attribute [C++17]}

The \texttt{[[fallthrough]]} attribute signals intentional fall-through in switch cases, silencing compiler warnings.
